\begin{apartats}

\apartat[1,25]{0,75}
Calcula les integrals indefinides $\displaystyle\int(x^3-2x+5)\,dx$ i
$\displaystyle\int\left(\frac1x-\frac1{x^2}\right)dx$.

\begin{solucio}
$\displaystyle\int(x^3-2x+5)\,dx=\frac{x^4}4-x^2+5x+k$ i
$\displaystyle\int\left(\frac1x-\frac1{x^2}\right)dx=\ln|x|+\frac1x+k$.
\end{solucio}

\begin{tria}{calcul-integral}
\itemtria{barrow}{1}{1,25}
Calcula $\displaystyle\int_1^e\left(\frac1x-\frac1{x^2}\right)dx$ i $\displaystyle\int_{-1}^1(x^3-2x+5)\,dx$.

\begin{solucio}
$\displaystyle\int_1^e\left(\frac1x-\frac1{x^2}\right)dx=\left[\ln x+\frac1x\right]_1^e
=\left(1+\frac1e\right)-(0+1)=\frac1e$.\\
$\displaystyle\int_{-1}^1(x^3-2x+5)\,dx=\left[\frac{x^4}4-x^2+5x\right]_{-1}^1
=\left(\frac14-1+5\right)-\left(\frac14-1-5\right)=10$.
\end{solucio}

\itemtria{valor-absolut}{1}{1,25}
Calcula $\displaystyle\int_{-1}^3|x-1|\,dx$.

\begin{solucio}
$|x-1|=1-x$ si $x\le1$ i $x-1$ si $x\ge1$:
\[
\int_{-1}^3|x-1|\,dx=\int_{-1}^1(1-x)\,dx+\int_1^3(x-1)\,dx
=\left[x-\frac{x^2}2\right]_{-1}^1+\left[\frac{x^2}2-x\right]_1^3=2+2=4.
\]
(Geomètricament, són dos triangles de base 2 i altura 2.)
\end{solucio}
\end{tria}

\begin{nomesllarg}
\apartat{0,75}
Calcula $\displaystyle\int_0^2(x^2+1)\,dx$, i fes-lo servir per calcular $\displaystyle\int_0^2(3x^2+3)\,dx$
sense cap primitiva més.

\begin{solucio}
$\displaystyle\int_0^2(x^2+1)\,dx=\left[\frac{x^3}3+x\right]_0^2=\frac83+2=\frac{14}3$. Com que
$3x^2+3=3(x^2+1)$, la segona val $3\cdot\frac{14}3=14$.
\end{solucio}
\end{nomesllarg}

\end{apartats}
